\section{HD-F1-T05: three unit translates and an explicit multiplier three}
\label{sec:f1-exact-division}

On the principal Eisenstein component of HD-F1-S07, every nonzero
rational Weil class $t$, with $\delta t$ integral and
$T=\int_A t^2>0$, has an actual integral signed representative
with multiplier $3\delta$ and effective parts satisfying
\[
 \boxed{c_t(A)\le
 \max\left\{3\delta,
       36\left\lfloor\frac{\delta\sqrt{2T}}3\right\rfloor\right\}
       \le17\delta\lceil\sqrt T\rceil.}
 \tag{F1T5.1}
\]
The construction uses three effective translates by the specified
Eisenstein unit. In particular it proves the integral inclusion
\[
 3\bigl(W_A\cap H^4(A,\mathbb Z)\bigr)
       \subseteq\operatorname{cl}\bigl(\operatorname{CH}^2(A)\bigr)
 \tag{F1T5.2}
\]
on this component. This is a consequence of the published Schoen
smoothing, the explicit quotient of HD-F1-S07, and the exact
maximal-order Weil-lattice calculation DEG\_L T6.24. It uses no
Markman, Severn or032 theorem premise.

\paragraph{1. An effective cycle with its entire class prescribed.}
The construction of HD-F1-S05--S07 supplies, on a nonempty open
part of the component, an effective integral surface $Z$ of
degree36 with
\[
 [Z]=\frac32\theta^2+w',\qquad
       w'\in W_A,\qquad \int_A(w')^2=18.
 \tag{F1T5.3}
\]
Here the smooth source $S$ has $[S]=\beta^2/6+w$.
HD-F1-S06 gives $K_S^2=16$ and $c_2(S)=8$, so the normal
Chern formula and self-intersection give $[S]^2=c_2(N)=8$.
Since $\beta^4=216$ and $\beta^2w=0$, squaring that class
gives $\int_{A_{\mathrm{source}}}w^2=8-216/36=2$.
For the actual degree-nine
quotient $q$, one has $q^*\theta=\beta$ and
$q_*[S]=[Z]$, the latter by the stabilizer/Euler-characteristic
calculation in HD-F1-S07. Thus
\[
 q_*\beta^2=9\theta^2,\qquad w'=q_*w,\qquad
 q^*q_*w=9w,
\]
and the projection formula gives
\[
 \int_A(w')^2
       =\frac19\int_{A_{\mathrm{source}}}(9w)^2=18.
 \tag{F1T5.4}
\]
These prove every coefficient in(F1T5.3).

The class $\theta^2/2$ is integral in cohomology. Indeed, in a
symplectic integral basis for a principal polarization,
$\theta=\sum_{i=1}^4x_i\wedge y_i$, so
\[
 \frac{\theta^2}2=\sum_{i<j}
                    x_i\wedge y_i\wedge x_j\wedge y_j.
 \tag{F1T5.5}
\]
Consequently $w'=[Z]-3\theta^2/2$ belongs to the actual integral
Weil lattice $\mathscr W_A$. Formula(F1T5.5) is used only for
integrality of a cohomology class; the signed cycles below are
constructed from $Z$ and its unit pullbacks.

The same full-class assertion holds at every fibre, with $Z$
allowed to be an effective cycle after specialization. To see this,
take the closure of a dominant component of the actual degree36
cycle family in the proper relative Chow space. Its image contains
a nonempty open set and is closed, hence is the whole irreducible
component. Its universal cycle class is flat locally. The relation
\[
 (1-e_W)\bigl([Z]-3\theta^2/2\bigr)=0
\]
and its square18 persist on that closure, because the projector,
polarization and cup pairing are flat. Degree remains36. This is
exactly the supplied proper Chow specialization argument, applied
to the full displayed class rather than just nonvanishing of its
projection. Thus(F1T5.3)--(F1T5.5) hold for an effective integral
cycle on each target fibre. A reducible special fibre causes no
change in the subsequent signed-cycle calculation.

\paragraph{2. The exact integral sublattice produced by unit differences.}
Write $\mathcal O=\mathbb Z[\omega]$, $\omega^2+\omega+1=0$,
and $\Omega=[\omega]^*$. The determinant coefficient action
on $\mathscr W_A$ is the one in DEG\_L T6.24; on the middle
fourfold Weil plane this geometric unit acts by
$\omega^4=\omega$. The lattice is a rank-one projective
$\mathcal O$-module. T6.24 gives the exact norm/index identity
\[
 [\mathscr W_A:\mathcal O w']
       =\frac{\int_A(w')^2}{2|\operatorname{disc}K|}=3.
 \tag{F1T5.6}
\]
There is a unique ideal of norm3 in $\mathcal O$, namely
$\mathfrak p=(1-\omega)$: modulo3 the polynomial
$X^2+X+1$ is $(X-1)^2$, and this describes the unique prime
above3. The index-three submodule is therefore
$\mathcal O w'=\mathfrak p\mathscr W_A$.
This can be checked after choosing a free $\mathcal O$ basis;
the existence of such a basis uses the Eisenstein class number one.

Set
\[
 u=(1-\Omega)[Z]=(1-\omega)w'.
 \tag{F1T5.7}
\]
Since $(1-\omega)^2=-3\omega$, equations(F1T5.6)--(F1T5.7)
give the exact equality of integral lattices
\[
 \mathcal O u=\mathfrak p^2\mathscr W_A=3\mathscr W_A.
 \tag{F1T5.8}
\]
Thus $u,\Omega u$ are a $\mathbb Z$ basis of $3\mathscr W_A$.
The positive middle pairing has
\[
 u^2=54,\qquad u\cdot\Omega u=-27,
                       \qquad(\Omega u)^2=54,
 \tag{F1T5.9}
\]
because multiplication by $1-\omega$ has norm3, and
$\operatorname{Re}\omega=-1/2$. This proves the exact lattice
comparison used to extract the class; no new choice of its
normalization has been made.

\paragraph{3. The actual integral effective parts and their degree.}
For $\delta t\in\mathscr W_A$, equation(F1T5.8) gives unique
integers $a,b$ with
\[
 3\delta t=a u+b\Omega u,
 \qquad a^2-ab+b^2=\frac{\delta^2T}{6}.
 \tag{F1T5.10}
\]
Let $Z_0=Z$, $Z_1=[\omega]^*Z$, $Z_2=[\omega^2]^*Z$.
These are actual effective integral cycles of degree36: the unit
is an automorphism preserving the polarization. The integral cycle
\[
 \Gamma=aZ_0+(b-a)Z_1-bZ_2
 \tag{F1T5.11}
\]
has class $3\delta t$. Split its displayed coefficients into
positive and negative parts. Their sums agree, and both equal
\[
 M=\frac{|a|+|b-a|+|b|}{2}
                =\max\{|a|,|b|,|a-b|\}.
 \tag{F1T5.12}
\]
For example the middle equality is the diameter of the three
real numbers $0,a,b$; ordering them proves the formula in all
sign cases. The resulting effective cycles $\Gamma^+,\Gamma^-$
each have degree $36M$. Their possible common components may be
retained, so the construction is valid without a disjoint-support
assumption.

The three completed-square identities for
$a^2-ab+b^2$ show
\[
 \frac34M^2\le a^2-ab+b^2=\frac{\delta^2T}{6},
 \qquad
 M\le\left\lfloor\frac{\delta\sqrt{2T}}3\right\rfloor.
 \tag{F1T5.13}
\]
For $a$ use $(b-a/2)^2+3a^2/4$, for $b$ interchange the
variables, and for $a-b$ use $((a+b)^2+3(a-b)^2)/4$.
This proves the first bound in(F1T5.1). Finally
$36M\le12\sqrt2\,\delta\sqrt T<17\delta\lceil\sqrt T\rceil$
and $3\delta\le17\delta\lceil\sqrt T\rceil$.
The claimed uniform bound follows, with the stated integral
multiplier and effective positive and negative representatives.

\paragraph{The four requested frameworks and row identifiers.}
For DEG\_L Q-DIV and Q-BOUNDS, the relative cycle family and its
proper Chow closure prove(F1T5.3) on every fibre; this is the
specific variational realization of Grothendieck1966, p.103,
footnote13, in the Hilbert/Chow setting of FGA221,
Theorems3.1--3.2 and the supplied Chow lemma. The polarization
biextension identity is $\widehat q\lambda_\theta q=\lambda_\beta$,
and the unit identities are
$\widehat\omega\lambda_\theta\omega=\lambda_\theta$;
these are SGA7 I,VII2.9.5--2.9.6 and VIII3.2 with the maps
of HD-F1-F03. The pure middle cohomology class represented by
(F1T5.11) is multiplied by $r^2$ under the degree-two Adams
character realization, SGA6 V,Appendix7.1--7.5. The corresponding
integral operation on $K_0$ and the trace/object data are
related by the forgetful morphism with kernel example(F1F4.5).
The finite-presentation and Betti/de Rham/$\ell$-adic realizations
are the ones of HD-F1-F04. In the archive frame, the actual cycle
(F1T5.11) supplies conditionA's integral representative; the
positive Gram(F1T5.9) is the conditionI calculation, and
$1-\Omega$ is an actual integral graph-correspondence polynomial.
These record the morphisms for the refined integral inclusion
(F1T5.2), without an additional conjectural premise.

The same degree-thirty-six input is also obtained directly from
the original canonical norm-image surface by HD-F1-S10, paragraph 7.
Its actual degree-nine quotient has generic degree one on the
surface, and supplies the same full class and integral Weil norm
eighteen used throughout this extraction.
